\makeatletter\let\ifGm@compatii\relax\makeatother
\documentclass[red]{beamer}

%\usepackage{beamerthemesplit}
\usepackage{listings}
\usepackage{comment}
\usepackage{url}
\usepackage{bm}
\usepackage{amsmath}
\usepackage{tikz}
\usetikzlibrary{arrows.meta}
\usetikzlibrary{shapes.symbols}

\usepackage{hyperref} % has to be last
\hypersetup{
    colorlinks=true,
    linkcolor=blue,
    filecolor=magenta,      
    urlcolor=cyan,
    pdfpagemode=FullScreen,
    }

\tikzset{
    saiarrow/.style={{-Latex}, thick}
}

\newcommand{\isred}[1]{{\color{red} {\bf #1}}} 
\newcommand{\isblue}[1]{{\color{blue}{\bf #1}}} 
\newcommand{\isgreen}[1]{{\color{green} {#1}}} 
\newcommand{\isdarkgreen}[1]{{\color{darkgreen} {#1}}} 
\newcommand{\isorange}[1]{{\color{orange} {#1}}} 

\newcommand{\N}{\mathbb N}
\newcommand{\Z}{\mathbb Z}
\newcommand{\Q}{\mathbb Q}
\newcommand{\R}{\mathbb R}

\begin{document}


\setbeamercolor{alerted_text}{fg=cyan}
\xdefinecolor{purplish}{cmyk}{0.75,0.75, 0, 0}
\xdefinecolor{purple}{cmyk}{0.75,0.75, 0, 0}
\colorlet{newred}{red!60!black}

\title{Session 6: CPSA Goals}
\author{Slides prepared by UMBC Protocol Analysis Lab}
\date{Last edited \today}
\begin{frame}
\maketitle
\end{frame}

\begin{frame}\frametitle{What is a security goal?}
\pause
A \emph{security goal} is a goal that we assign for a given protocol to evaluate if it is sufficiently secure for application.

\pause
Security goals formalize our notion of ``secure''.

\bigskip
\pause
Not every protocol has the same requirement for security.

Common Security Goals
\begin{itemize}
\pause
\item \textbf{Confidentiality:} In a two-party communication protocol, confidentiality holds if no third party is able to access the information sent.

\pause
\item \textbf{Agreement:} In a two-party communication protocol, agreement holds if both parties agree on some set of data values in the protocol.
These values are then used to derive the key.
\end{itemize}

\end{frame}

\begin{frame}\frametitle{Other common security goals}
\begin{enumerate}
\item Aliveness
\item Weak Agreement
\item Non-injective Agreement
\item A specific relay party cannot read the message
\item The message must be timestamped
\item Perfect forward secrecy
\item Crypto primitives are resistant to a predefined attack
\end{enumerate}
\end{frame}

\begin{frame}\frametitle{Common protocol security goals}
\pause
\textbf{Authentication protocols:}
\pause

Want to ensure that the other person is who they say they are.
What security goals?
\pause

Agreement needed.

\pause
\textbf{Network protocols:}
\pause

Need to send a lot of data quickly.
What security goals?
\pause

Confidentiality needed.

Note that the ``second'' party need not have agreement (in case packets dropped).

\pause
\textbf{Key Exchange protocols:}
\pause

Want to distribute keys to the appropriate participants on the network.
What security goals?
\pause

Confidentiality and Agreement (on the secret) needed.


\end{frame}


\begin{frame}\frametitle{CPSA Goals}
\begin{itemize}
\item Protocol goals as mathematical statements
\item Find strands that satisfy (or don't satisfy) some set of properties
\item May reduce time taken to terminate
\end{itemize}
\end{frame}

\begin{frame}\frametitle{Some more tips for reading specifications}

\begin{itemize}
\item FIND THE SECURITY GOALS
\begin{itemize}
\item Are additional SGs needed?
\end{itemize}
\item Modularize the spec into blocks of messages
\begin{itemize}
\item ``In the first set of exchanges''
\item ``[party1] preshares [data] with [party2]''
\end{itemize}
\item Ignore nonsense
\begin{itemize}
\item ``[protocol] does not consider [niche usecase] because\dots''
\item ``We assume clients have access to an ALU'' (structural vs.\ implementation)
\item ``Our conformance program handles...''
\end{itemize}
\end{itemize}
\end{frame}

\begin{frame}\frametitle{Needham-Schroeder}
\begin{itemize}
\item Exactly the same as before
\pause
\item We know about the fix, so we will check that too.
\pause 
\item Some defined security goals that we can model:
\begin{enumerate}
\pause
\item (NS) Agree on name of Bob
\pause
\item (NSL) None of the nonces are revealed
\pause
\item (NSL) Do we get one-way authentication?
\end{enumerate}
\pause
\item Let's try modeling it!
\end{itemize}
\end{frame}

\begin{frame}\frametitle{Conclusion}
\begin{itemize}
\pause
\item Does Needham-Schroeder work?
\pause
\begin{itemize}
\item What does ``work'' mean?
\end{itemize}
\pause
\item What about NSL?
\pause
\end{itemize}
\end{frame}



\begin{frame}\frametitle{Survey}
Fill in details about the survey here!
\end{frame}
\end{document}